\section{Continuous Assessment 2 (CA-2)}

\ExamHeader{Continuous Assessment 2 (CA-2)}{Project Portfolio Set D}{Continuous Submission}{30 Marks}{None}

\noindent
\textbf{Continuous Assessment 2 (CA-2) Project Assignment Guidelines:}
\begin{enumerate}[topsep=2pt,itemsep=1pt]
    \item CA-2 evaluates in-depth applied physics project investigations, analytical derivations, and engineering design reports.
    \item Each report must include: \textbf{Theoretical Foundation}, \textbf{Mathematical Model / Derivation}, \textbf{Engineering Application Case Study}, and \textbf{Performance Analysis}.
\end{enumerate}
\vspace{3mm}

\begin{projectbox}{1: Electrodynamic Analysis of Maxwell's Equations and Waveguide Propagation}
\textbf{Project Overview:} Rigorous investigation of electromagnetic wave propagation in rectangular waveguides and displacement current verification in dielectric media.

\textbf{1. Theoretical Background:}
Maxwell's equations synthesize electrostatics, magnetostatics, and time-varying electrodynamics. The displacement current term $\frac{\partial \vec{D}}{\partial t}$ resolves the charge conservation paradox in capacitor charging circuits.

\textbf{2. Mathematical Formulation \& Governing Equations:}
In free space, $\nabla \times (\nabla \times \vec{E}) = -\frac{\partial}{\partial t}(\nabla \times \vec{B}) = -\mu_0 \epsilon_0 \frac{\partial^2 \vec{E}}{\partial t^2}$.
Using $\nabla \times (\nabla \times \vec{E}) = \nabla(\nabla \cdot \vec{E}) - \nabla^2 \vec{E} = -\nabla^2 \vec{E}$, we obtain:
\begin{equation}
\nabla^2 \vec{E} = \frac{1}{c^2}\frac{\partial^2 \vec{E}}{\partial t^2}, \quad c = \frac{1}{\sqrt{\mu_0 \epsilon_0}} \approx 3 \times 10^8\text{ m/s}
\end{equation}

\textbf{3. Engineering Case Study \& Practical Implementation:}
Design of a TE10 mode microwave waveguide operating at $10\text{ GHz}$ for radar telecommunications.

\textbf{4. Critical Findings \& Performance Analysis:}
Waveguide cut-off frequency $f_c = \frac{c}{2a}$ strictly limits propagation. Group velocity $v_g = c\sqrt{1 - (f_c/f)^2}$ remains subluminal, preserving relativistic causality.
\end{projectbox}

\vspace{4mm}

\begin{projectbox}{2: Quantum Mechanical Modeling of 1D Potential Wells and Tunneling Dynamics}
\textbf{Project Overview:} Numerical and analytical investigation of electron confinement in semiconductor quantum wells and scanning tunneling microscopy (STM).

\textbf{1. Theoretical Background:}
When a particle is confined in an infinite potential box of length $L$, boundary conditions quantize momentum and energy into discrete eigenvalues.

\textbf{2. Mathematical Formulation \& Governing Equations:}
Eigenvalues and eigenfunctions:
\begin{equation}
E_n = \frac{n^2 \pi^2 \hbar^2}{2m L^2}, \quad \psi_n(x) = \sqrt{\frac{2}{L}}\sin\left(\frac{n\pi x}{L}\right)
\end{equation}
For finite barriers of height $V_0 > E$, the transmission probability is $T \approx e^{-2\alpha L}$ where $\alpha = \frac{\sqrt{2m(V_0 - E)}}{\hbar}$.

\textbf{3. Engineering Case Study \& Practical Implementation:}
Analysis of GaAs/AlGaAs quantum-well lasers and atomic-resolution imaging in Scanning Tunneling Microscopes.

\textbf{4. Critical Findings \& Performance Analysis:}
Exponential dependence of tunnel current $I \propto e^{-2\alpha d}$ yields sub-angstrom vertical resolution for surface topographical mapping.
\end{projectbox}

\vspace{4mm}

